\chapter{Markov Chains}

Up to now, our random sequences were almost always built from \emph{independent}
building blocks: repeated coin tosses, dice rolls, or measurement noise.
In every step, the next random variable knew nothing about the ones that came before.
Independence makes calculations clean, but the world rarely obliges. Tomorrow's weather
depends on today's. Your next word depends on the word you just typed. Whether a gambler is
up or down at the next bet depends on whether they are up or down right now. The genome you
inherit at one site is correlated with the site next door.

We need a model of randomness \emph{with memory} -- but with structured dependence so
that exact calculations remain tractable. The Markov chain provides this framework.
It is the simplest tractable model in which the future depends on the past, and it is widely applied
in probability and statistics: web page ranking, language modeling, ion channels in biophysics,
queueing theory, Markov chain Monte Carlo in Bayesian computation, and the spread of
epidemics can all be formulated as Markov chains.

The core property of a Markov chain is conditional independence: \emph{the present screens
off the past}. Conditioned on the state of the process at the current time, its future trajectory
is independent of its past history. This chapter develops the mathematical framework that
flows from this single property.

\chapterline

\section{The Markov property: the present is enough}
\label{ch23:markov-property}

Picture a system hopping among a set of states (such as a particle moving between discrete locations).
At each step, the transition to the next state depends only on the state it occupies
\emph{right now}. The transition rule ignores the path taken to reach the current state;
the current location alone determines the transition probabilities.

\begin{intuitionbox}
A Markov chain is a process hopping among a set of \emph{states} where the only thing that matters
for the next step is the current state. The history is not erased -- you can still write it
down -- but it is \emph{irrelevant} to the future once the present is known. ``Given the
present, the future and the past are independent.''
\end{intuitionbox}

Let us make this precise. We have a sequence of random variables $X_0, X_1, X_2, \dots$,
each taking values in a finite or countable set $S$ called the \emph{state space} -- the catalog
of all possible locations, conditions, or configurations the system can inhabit. We think
of $X_n$ as the state at time $n$.

\begin{definition}[Markov property]
\label{ch23:def-markov}
The sequence $(X_n)_{n \ge 0}$ is a (discrete-time) \emph{Markov chain} on state space $S$
if for every $n \ge 0$ and all states $i_0, i_1, \dots, i_{n-1}, i, j \in S$ with positive
probability of the conditioning event,
\[
  \P(X_{n+1} = j \given X_n = i,\, X_{n-1} = i_{n-1}, \dots, X_0 = i_0)
   = \P(X_{n+1} = j \given X_n = i).
\]
The chain is \emph{time-homogeneous} if this last quantity does not depend on $n$ -- meaning
the transition rules stay identical as the clock ticks forward. We then write
$p_{ij} = \P(X_{n+1} = j \given X_n = i)$ for the \emph{transition probability} from $i$ to $j$,
the probability of hopping directly from $i$ to $j$ in a single tick. Unless we say otherwise,
all our chains are time-homogeneous.
\end{definition}

Read the equation slowly. On the left you condition on the \emph{entire} history. On the
right you keep only the most recent value $X_n = i$. The Markov property says these two are
equal: the extra history $i_{n-1}, \dots, i_0$ adds nothing once you know $i$. That is the
``memoryless given the present'' slogan made into mathematics.

\begin{trapbox}
``Memoryless'' does \emph{not} mean the steps are independent, and it does \emph{not} mean
the chain forgets everything. The system's current state is highly informative about its next
state. What the Markov property says is narrower: \emph{additional} past beyond the present
contributes no \emph{further} information. Also beware the symmetric trap -- the past and
future are conditionally independent \emph{given the present}, but they are usually very much
dependent if you do not condition. Knowing $X_0$ tells you plenty about $X_5$.
\end{trapbox}

\begin{whybox}
Why is this the ``right'' weakening of independence? Independence says
$\P(X_{n+1}=j \given \text{past}) = \P(X_{n+1}=j)$ -- the past is irrelevant, full stop.
That is too strong for memory. The opposite extreme, letting $X_{n+1}$ depend on the whole
history, is too rich: to specify the model you would need a conditional probability for every
possible history, and histories grow without bound. The Markov property is the unique
middle ground where (a) the future genuinely depends on the past, yet (b) the dependence is
funneled entirely through a single current state, so the model is described by a finite
table. It is the most memory you can have while staying finitely specifiable.
\end{whybox}

\begin{historybox}
The chains are named for \textbf{Andrey Andreyevich Markov} (1856--1922), a student of
Pafnuty Chebyshev at Saint Petersburg University. In 1906, Markov was spurred into action by
an ideological controversy. Pavel Nekrasov, a mathematician and theologian, had asserted that
the weak law of large numbers strictly required independence, enlisting probability to argue
that voluntary human decisions must be mutually independent acts of free will. Markov, an
uncompromising materialist with little patience for theological mathematics, set out to
shatter Nekrasov's premise. He asked: could limit theorems still hold when variables are
entangled with their neighbors?

Markov had no modern measure theory -- Andrey Kolmogorov's foundational axioms would not appear
until 1933 -- and no graph-theoretic language of networks. Working with classical tools like
determinants, continued fractions, and Chebyshev's inequality, Markov had to invent matrix
contraction techniques from scratch. To prove that his dependent variables obeyed a law of
large numbers, he introduced what we now call the Markov property: the future is conditioned
only on the immediate present. In 1913, to demonstrate his theory on empirical data, Markov
performed a remarkably modern experiment. He took the first 20{,}000 letters of Alexander
Pushkin's verse novel \emph{Eugene Onegin}, tallied vowels and consonants by hand, and showed
that a vowel was followed by another vowel with probability $0.128$, but by a consonant with
probability $0.872$. That two-state chain was the world's first statistical language model --
the direct intellectual ancestor of $n$-gram models and modern generative language models.

Markov's breakthrough unlocked modern stochastic processes: the Chapman--Kolmogorov equations
in physics, queueing theory in telecommunications, and Markov Chain Monte Carlo (MCMC)
algorithms that power computational Bayesian statistics and machine learning. Today, active
frontiers remain wide open. Chief among them is bounding \emph{mixing times} in high
dimensions -- understanding the sharp ``cutoff phenomenon'' where a chain hovers far from
equilibrium before abruptly dropping into steady state (as when riffle-shuffling cards), and
designing rapid-mixing non-reversible chains for sampling non-convex energy landscapes.
\end{historybox}

\subsection{From the present to the whole future}

The Markov property looks like a statement about one step, but it quietly controls the entire
trajectory. The probability of any specific path factors into a product of one-step pieces.

\begin{proposition}[Path probabilities multiply]
\label{ch23:path-factorization}
For a Markov chain with initial distribution $\P(X_0 = i_0) = \mu_{i_0}$,
\[
  \P(X_0 = i_0,\, X_1 = i_1,\, \dots,\, X_n = i_n)
   = \mu_{i_0}\, p_{i_0 i_1}\, p_{i_1 i_2}\, \cdots\, p_{i_{n-1} i_n}.
\]
\end{proposition}

\begin{whybox}
Use the chain rule of probability, which always holds:
\[
  \P(X_0=i_0,\dots,X_n=i_n) = \P(X_0=i_0)\prod_{k=1}^{n}\P(X_k=i_k \given X_0=i_0,\dots,X_{k-1}=i_{k-1}).
\]
Now apply the Markov property to each factor: the conditioning collapses to the immediately
preceding state, $\P(X_k=i_k \given X_{k-1}=i_{k-1}) = p_{i_{k-1}i_k}$. The general chain
rule is unconditional truth; the Markov property is what lets us throw away all but the last
condition. The trajectory's probability is just a walk along a sequence of edges, each
contributing its own factor -- start weight times a product of transition weights.
\end{whybox}

This factorization is the computational engine of the subject. A path is a product, and sums
over paths are exactly what matrix multiplication computes -- the bridge we build next.

\begin{recoverybox}
\textbf{Pause and think: identifying the Markov boundary.}
Before moving on to matrices, test whether your intuition can distinguish what is Markovian from what is not:
\begin{enumerate}[label=(\alph*)]
  \item You roll a fair die repeatedly. Let $S_n = X_1 + \dots + X_n$ be the running sum after $n$ rolls. Is $(S_n)_{n \ge 1}$ a Markov chain?
  \item Let $M_n = \max(X_1, \dots, X_n)$ be the running maximum roll. Is $(M_n)_{n \ge 1}$ a Markov chain?
  \item Let $Y_n = \ind(X_n > X_{n-1})$ record whether the $n$-th roll strictly exceeded the previous one. Is $(Y_n)_{n \ge 2}$ a Markov chain?
\end{enumerate}
\emph{Reflections:} In (a), yes: $S_{n+1} = S_n + X_{n+1}$. Given the current sum $S_n$, the next sum depends only on the fresh roll $X_{n+1}$; the breakdown of how $S_n$ was formed is irrelevant. In (b), yes: $M_{n+1} = \max(M_n, X_{n+1})$, which depends only on the current maximum and the new roll. But in (c), \textbf{no}: knowing only that $Y_n = 1$ tells you the roll went up, but not \emph{what} number was rolled. If $X_n = 6$, $Y_{n+1}$ is guaranteed to be $0$; if $X_n = 2$, $Y_{n+1}$ is likely $1$. To make the process Markovian, the state must store the actual value of $X_n$, not merely the binary direction of change.
\end{recoverybox}

\section{Transition matrices and state diagrams}
\label{ch23:transition-matrix}

Collect all the one-step probabilities into a square array. If $S = \{1, 2, \dots, m\}$, the
\emph{transition matrix} is
\[
  P = \big(p_{ij}\big)_{i,j \in S},
  \qquad p_{ij} = \P(X_{n+1} = j \given X_n = i).
\]
Row $i$ of $P$ is the probability distribution of where you go \emph{from} state $i$. Because
from $i$ you must land \emph{somewhere}, every row sums to $1$:
\[
  \sum_{j \in S} p_{ij} = 1 \quad \text{for every } i, \qquad p_{ij} \ge 0.
\]
A nonnegative square matrix whose rows each sum to $1$ is called a \emph{stochastic matrix} --
a square grid of probabilities where each row records how a total probability mass of $1$
sitting at state $i$ is dispersed among all possible destinations. Stochastic matrices and
(time-homogeneous, finite) Markov chains are the same thing seen two ways.

\begin{trapbox}
\emph{Rows} sum to $1$, not columns. Mixing this up is the single most common Markov-chain
bug. The convention ``state-as-row, distribution-as-row'' also forces us to put probability
vectors on the \emph{left}: a row vector $\pi$ updates by $\pi \mapsto \pi P$, not $P\pi$.
Some books transpose everything and multiply on the right. Pick one convention and hold it.
We use rows throughout.
\end{trapbox}

A picture says it faster than a matrix. Draw one node per state and an arrow $i \to j$
labeled $p_{ij}$ for each positive transition probability. This is the \emph{state diagram}.

\begin{figure}[H]
\centering
\begin{tikzpicture}[->,>=Stealth,node distance=3.2cm,auto,
  mystate/.style={circle,draw=sourceblue,thick,minimum size=1.1cm}]
  \node[mystate] (S) {Sun};
  \node[mystate,right=of S] (R) {Rain};
  \path
    (S) edge[loop left]  node {$0.8$} (S)
    (S) edge[bend left=18] node[above] {$0.2$} (R)
    (R) edge[bend left=18] node[below] {$0.4$} (S)
    (R) edge[loop right] node {$0.6$} (R);
\end{tikzpicture}
\caption{A two-state weather chain. From a sunny day, tomorrow is sunny with probability
$0.8$ and rainy with probability $0.2$; from a rainy day, sunny with probability $0.4$ and
rainy with probability $0.6$. Every arrow leaving a node corresponds to one entry of that
node's row, and those arrows' labels sum to $1$.}
\label{ch23:fig-weather-diagram}
\end{figure}

In matrix form, with the order (Sun, Rain),
\[
  P = \begin{pmatrix} 0.8 & 0.2 \\ 0.4 & 0.6 \end{pmatrix}.
\]
Check the rows: $0.8 + 0.2 = 1$ and $0.4 + 0.6 = 1$. Good.

\begin{example}[A reader who skims]
\label{ch23:ex-skim}
Model your reading mood with three states: Focused ($F$), Skimming ($K$), Distracted ($D$).
Each paragraph you transition according to
\[
  P = \begin{pmatrix}
   0.7 & 0.2 & 0.1 \\
   0.3 & 0.5 & 0.2 \\
   0.2 & 0.3 & 0.5
  \end{pmatrix}
  \quad
  \begin{array}{l} F \\ K \\ D \end{array}
\]
where rows and columns are ordered $F, K, D$. Reading off row $F$: if you are focused now,
next paragraph you stay focused with probability $0.7$, start skimming with $0.2$, drift off
with $0.1$. Figure~\ref{ch23:fig-skim-diagram} shows the complete state transition network.
Every state can reach every other state in a single step, forming an irreducible flow network.
\end{example}

\begin{figure}[H]
\centering
\begin{tikzpicture}[->,>=Stealth,
  mystate/.style={circle,draw=#1,thick,fill=#1!12,minimum size=1.15cm,font=\small\bfseries}]
  % Nodes arranged in an equilateral triangle
  \node[mystate=sourceblue] (F) at (0,2.6) {$F$};
  \node[mystate=quotientorange] (K) at (-2.4,-1.0) {$K$};
  \node[mystate=trapred] (D) at (2.4,-1.0) {$D$};

  % Self loops
  \path (F) edge[loop above, thick, looseness=4] node[above, font=\footnotesize] {$0.7$} (F);
  \path (K) edge[loop left, thick, looseness=4] node[left, font=\footnotesize] {$0.5$} (K);
  \path (D) edge[loop right, thick, looseness=4] node[right, font=\footnotesize] {$0.5$} (D);

  % Directed edges between F and K
  \path (F) edge[bend right=18, thick, sourceblue] node[left=2pt, font=\footnotesize] {$0.2$} (K);
  \path (K) edge[bend right=18, thick, quotientorange] node[right=2pt, font=\footnotesize] {$0.3$} (F);

  % Directed edges between F and D
  \path (F) edge[bend left=18, thick, sourceblue] node[right=2pt, font=\footnotesize] {$0.1$} (D);
  \path (D) edge[bend left=18, thick, trapred] node[left=2pt, font=\footnotesize] {$0.2$} (F);

  % Directed edges between K and D
  \path (K) edge[bend left=18, thick, quotientorange] node[above=1pt, font=\footnotesize] {$0.2$} (D);
  \path (D) edge[bend left=18, thick, trapred] node[below=1pt, font=\footnotesize] {$0.3$} (K);

\end{tikzpicture}
\caption{The three-state reading mood chain of Example~\ref{ch23:ex-skim}. From Focused ($F$), Skimming ($K$), or Distracted ($D$), the probability transitions along directed paths. The outgoing weights from any node sum to $1$. Because every state communicates directly with every other state, the chain is irreducible and aperiodic.}
\label{ch23:fig-skim-diagram}
\end{figure}

\begin{walkbox}
A particle does a random walk on the corners of a triangle, at each step moving to one of the
two \emph{other} corners with equal probability $1/2$, never staying put. Write its
transition matrix. (Answer: label the corners $1,2,3$; then $p_{ii}=0$ and
$p_{ij}=1/2$ for $i \ne j$. Each row is $(0,\tfrac12,\tfrac12)$ in some order, summing to
$1$.) Notice the zero diagonal -- the walk \emph{cannot} stay. One might ask whether this forces a kind of rhythm called periodicity -- a question we return to in Example~\ref{ch23:ex-triangle-period}.
\end{walkbox}

\section{$n$-step transitions are matrix powers}
\label{ch23:matrix-powers}

Now the central question. You know the one-step probabilities.
What is the probability of going from $i$ to $j$ in \emph{exactly $n$ steps}? Write this
$n$-step transition probability as
\[
  p_{ij}^{(n)} = \P(X_n = j \given X_0 = i).
\]

Start with two steps. To get from $i$ to $j$ in two steps you must pass through \emph{some}
intermediate state $k$ at time $1$. The events ``go through $k$'' are mutually exclusive and
exhaustive, so we add them up:
\[
  p_{ij}^{(2)} = \sum_{k \in S} \P(X_1 = k \given X_0 = i)\,\P(X_2 = j \given X_1 = k)
              = \sum_{k \in S} p_{ik}\, p_{kj}.
\]
The middle step used the Markov property: once we are at $k$ at time $1$, the second jump
ignores that we came from $i$. But $\sum_k p_{ik}\,p_{kj}$ is exactly the $(i,j)$ entry of
the matrix product $P \cdot P = P^2$. The two-step probabilities \emph{are} the entries of
$P^2$.

\begin{intuitionbox}
Matrix multiplication and ``sum over all intermediate states'' are literally the same
operation. The row-times-column rule $\sum_k p_{ik} p_{kj}$ enumerates every two-step path
$i \to k \to j$, multiplies the probabilities along each path (paths are products), and adds
across paths (alternatives add). Probability bookkeeping \emph{is} linear algebra.
\end{intuitionbox}

The same argument with $n$ split as $r + s$ gives the general law.

\begin{theorem}[Chapman--Kolmogorov]
\label{ch23:chapman-kolmogorov}
For all $r, s \ge 0$,
\[
  p_{ij}^{(r+s)} = \sum_{k \in S} p_{ik}^{(r)}\, p_{kj}^{(s)},
  \qquad\text{equivalently}\qquad
  P^{(r+s)} = P^{(r)} P^{(s)}.
\]
Consequently $p_{ij}^{(n)} = (P^n)_{ij}$: the $n$-step transition matrix is the $n$-th power
of $P$.
\end{theorem}

\begin{whybox}
Condition on where the chain is at the intermediate time $r$. To travel from $i$ to $j$ in
$r+s$ steps, it sits at some state $k$ at time $r$; the first leg ($i \to k$ in $r$ steps) and
the second leg ($k \to j$ in $s$ steps) are governed by $P^{(r)}$ and $P^{(s)}$, and the
Markov property makes the second leg independent of how the first was traveled. Summing over
$k$ is the law of total probability. Writing it as matrices, ``sum over the middle index'' is
the definition of the product, so $P^{(r+s)} = P^{(r)}P^{(s)}$. Setting $r=s=1$ gives
$P^{(2)} = P^2$, and induction gives $P^{(n)} = P^n$.
\end{whybox}

\begin{connectionbox}
This is where the linear algebra book pays off. The $n$-step behavior of a chain is the $n$-th
power of a matrix, and powers of matrices are understood through \emph{eigenvalues}. If $P$ is
diagonalizable as $P = V \Lambda V^{-1}$ with eigenvalues $\lambda_1, \dots, \lambda_m$ on the
diagonal of $\Lambda$, then $P^n = V \Lambda^n V^{-1}$, and $\Lambda^n$ just raises each
eigenvalue to the $n$. A stochastic matrix always has $\lambda_1 = 1$ as an eigenvalue (the
all-ones vector $\mathbf{1}$ satisfies $P\mathbf{1} = \mathbf{1}$, since rows sum to $1$), and
all eigenvalues satisfy $\abs{\lambda} \le 1$. As $n \to \infty$, every $\abs{\lambda} < 1$
term decays like $\lambda^n \to 0$, and only the $\lambda = 1$ part survives. \emph{That
surviving piece is the long-run equilibrium of the chain.} The whole theory of convergence we
meet in the next chapter is, underneath, the statement ``the second-largest eigenvalue
controls the rate.''
\end{connectionbox}

\subsection{Pushing a distribution forward}

If you are uncertain about the starting state, encode that uncertainty as a row vector
$\mu^{(0)} = (\mu^{(0)}_1, \dots, \mu^{(0)}_m)$ with $\mu^{(0)}_i = \P(X_0 = i)$. The
distribution after one step is obtained by the law of total probability,
\[
  \mu^{(1)}_j = \P(X_1 = j) = \sum_i \P(X_0 = i)\,p_{ij} = \big(\mu^{(0)} P\big)_j,
\]
so $\mu^{(1)} = \mu^{(0)} P$, and by induction $\mu^{(n)} = \mu^{(0)} P^n$. A distribution is a
row vector; time advances by right-multiplication by $P$.

\begin{example}[Two steps of the weather chain]
\label{ch23:ex-weather-twostep}
Take $P = \left(\begin{smallmatrix} 0.8 & 0.2 \\ 0.4 & 0.6 \end{smallmatrix}\right)$. Then
\[
  P^2 = \begin{pmatrix} 0.8 & 0.2 \\ 0.4 & 0.6 \end{pmatrix}
        \begin{pmatrix} 0.8 & 0.2 \\ 0.4 & 0.6 \end{pmatrix}
      = \begin{pmatrix} 0.72 & 0.28 \\ 0.56 & 0.44 \end{pmatrix}.
\]
So if today is sunny, the probability it is sunny the day \emph{after} tomorrow is
$p_{SS}^{(2)} = 0.72$. Check the arithmetic of the top-left entry the path way: $S \to S \to
S$ contributes $0.8 \times 0.8 = 0.64$, and $S \to R \to S$ contributes $0.2 \times 0.4 =
0.08$; together $0.64 + 0.08 = 0.72$. Two-step probability, two paths, products added.
Suppose instead today's weather is itself uncertain, $\mu^{(0)} = (0.5, 0.5)$. Then
$\mu^{(1)} = (0.5,0.5)P = (0.6, 0.4)$ and $\mu^{(2)} = (0.6,0.4)P = (0.64, 0.36)$.
\end{example}

\begin{intuitionbox}
\textbf{The dual view: fluid and pebble flows.}
There are two complementary ways to think about a Markov chain:
\begin{enumerate}
  \item \textbf{Microscopic (the solitary walker):} A single token jumps probabilistically from node to node, choosing its next destination by rolling a die with probabilities $p_{ij}$.
  \item \textbf{Macroscopic (the flow network):} Imagine pouring $1{,}000$ pebbles (or a continuous volume of fluid) onto the network. At each tick, the pebbles currently resting on node $i$ divide up along outgoing pipes: fraction $p_{ii}$ stays in the reservoir, while fraction $p_{ij}$ flows through the pipe to node $j$.
\end{enumerate}
Under this macroscopic view, probability distribution updates $\mu^{(t+1)} = \mu^{(t)}P$ are simply the physical redistribution of pebbles. When the incoming flow along pipe $j \to i$ balances the outgoing flow along pipe $i \to j$, the system reaches a steady-state fluid equilibrium.
\end{intuitionbox}

\begin{figure}[H]
\centering
\begin{tikzpicture}[
  tank/.style={circle, draw=#1, thick, fill=#1!15, minimum size=1.05cm, font=\small\bfseries, align=center},
  flow/.style={->, >=Stealth, thick}
]
  % Lane 0: t=0
  \node[tank=sourceblue] (S0) at (0, 3.6) {$S$\\ \scriptsize $1000$};
  \node[tank=targetgreen] (R0) at (3.0, 3.6) {$R$\\ \scriptsize $0$};
  \node[font=\footnotesize\bfseries, left=1.15cm of S0] {Step $t=0$};
  \node[font=\scriptsize, gray, align=center] at (1.5, 4.3) {all mass on $S$};

  % Lane 1: t=1
  \node[tank=sourceblue] (S1) at (0, 1.8) {$S$\\ \scriptsize $800$};
  \node[tank=targetgreen] (R1) at (3.0, 1.8) {$R$\\ \scriptsize $200$};
  \node[font=\footnotesize\bfseries, left=1.15cm of S1] {Step $t=1$};
  \draw[flow, sourceblue] (S1) edge[bend left=16] node[above, font=\scriptsize] {$200 \to$} (R1);
  \draw[flow, sourceblue] (S1) edge[loop left, looseness=4] node[left, font=\scriptsize] {$800$} (S1);

  % Lane 2: t=2
  \node[tank=sourceblue] (S2) at (0, 0.0) {$S$\\ \scriptsize $720$};
  \node[tank=targetgreen] (R2) at (3.0, 0.0) {$R$\\ \scriptsize $280$};
  \node[font=\footnotesize\bfseries, left=1.15cm of S2] {Step $t=2$};
  \draw[flow, sourceblue] (S2) edge[bend left=16] node[above, font=\scriptsize] {$160 \to$} (R2);
  \draw[flow, targetgreen] (R2) edge[bend left=16] node[below, font=\scriptsize] {$\gets 80$} (S2);

  % Lane 3: Steady state
  \node[tank=sourceblue] (Sinf) at (7.4, 1.8) {$S$\\ \scriptsize $667$};
  \node[tank=targetgreen] (Rinf) at (10.4, 1.8) {$R$\\ \scriptsize $333$};
  \node[font=\footnotesize\bfseries, above=0.85cm of $(Sinf)!0.5!(Rinf)$] {Steady State $t \to \infty$};
  \draw[flow, sourceblue] (Sinf) edge[bend left=20] node[above, font=\scriptsize] {$133.3 \to$} (Rinf);
  \draw[flow, targetgreen] (Rinf) edge[bend left=20] node[below, font=\scriptsize] {$\gets 133.3$} (Sinf);
  \node[font=\scriptsize, align=center, below=0.85cm of $(Sinf)!0.5!(Rinf)$] {
    \textbf{Dynamic Balance:}\\
    flow $S \to R = 667 \times 0.2 = 133.3$\\
    flow $R \to S = 333 \times 0.4 = 133.3$
  };

  % Divider
  \draw[dashed, gray!50, thick] (4.8, -0.6) -- (4.8, 4.5);

\end{tikzpicture}
\caption{The pebble distribution flow across the weather chain. Starting with $1{,}000$ pebbles at state Sun ($t=0$), each tick redistributes pebbles across the directed pipes. By $t=2$, Sun holds $720$ and Rain holds $280$. In the infinite limit, the distribution settles at $(\tfrac23, \tfrac13)$: exactly $133.3$ pebbles cross in each direction every step, maintaining constant reservoir populations.}
\label{ch23:fig-pebble-flow}
\end{figure}

Watch a chain forget where it started. Iterate $\mu^{(n)} = \mu^{(0)}P^n$ for the
weather chain from two very different starting beliefs, and the probability of ``Sun'' converges
to the same value no matter where we began.

\begin{figure}[H]
\centering
\begin{tikzpicture}
\begin{axis}[
  width=12cm, height=6.2cm,
  xlabel={day $n$}, ylabel={$\P(X_n = \text{Sun})$},
  xmin=0, xmax=10, ymin=0, ymax=1.05,
  xtick={0,2,4,6,8,10}, ytick={0,0.3333,0.6667,1},
  yticklabels={$0$,$\tfrac13$,$\tfrac23$,$1$},
  title style={font=\small},
  title={Convergence of marginal distribution to stationary state $\pi_{\text{Sun}} = 2/3$},
  legend style={font=\small, at={(0.98,0.06)}, anchor=south east, column sep=4pt},
  legend cell align=left,
  grid=both, grid style={gray!18},
  every axis plot/.append style={thick}]
  % start sunny: P(Sun)=1; recursion s_{n+1}=0.8 s + 0.4(1-s)=0.4+0.4 s
  \addplot[sourceblue,mark=*,mark size=1.4pt] coordinates {
    (0,1)(1,0.8)(2,0.72)(3,0.688)(4,0.6752)(5,0.67008)(6,0.668032)(7,0.6672128)(8,0.66688512)(9,0.666754048)(10,0.6667016192)};
  \addlegendentry{start Sun}
  % start rainy: P(Sun)=0
  \addplot[densityred,mark=square*,mark size=1.3pt] coordinates {
    (0,0)(1,0.4)(2,0.56)(3,0.624)(4,0.6496)(5,0.65984)(6,0.663936)(7,0.6655744)(8,0.66622976)(9,0.666491904)(10,0.6665967616)};
  \addlegendentry{start Rain}
  \addplot[gray,dashed,domain=0:10,samples=2] {0.6667};
  \addlegendentry{$\pi_{\text{Sun}} = 2/3$}
\end{axis}
\end{tikzpicture}
\caption{The weather chain forgets its start. Whether day $0$ is certainly sunny (blue) or
certainly rainy (red), the probability of sun on day $n$ rushes to the same limit
$2/3$. The gap between the two curves shrinks geometrically -- it is multiplied by the
second eigenvalue $\lambda_2 = 0.4$ each day. This convergence to a unique equilibrium is the
subject of the next chapter; here it is a teaser for what matrix powers do.}
\label{ch23:fig-weather-converge}
\end{figure}

\begin{happeningbox}
Why do both curves land on $2/3$? The eigenvalues of the weather $P$ are $\lambda_1 = 1$ and
$\lambda_2 = 0.8 + 0.6 - 1 = 0.4$ (for a $2\times2$ stochastic matrix the second eigenvalue is
always $\operatorname{tr}P - 1$). Writing $\mu^{(0)} = c_1 v_1 + c_2 v_2$ in the eigenbasis,
$\mu^{(n)} = c_1 v_1 + c_2 (0.4)^n v_2$. The $\lambda=1$ part $c_1 v_1$ is fixed; the $0.4$
part decays to nothing. What is left is the equilibrium row vector, the left eigenvector for
$\lambda = 1$, which works out to $(2/3, 1/3)$. The starting distribution only enters through
the dying term, so it is forgotten. The ``memory'' of the initial condition has a half-life of
about $\log(1/2)/\log(0.4) \approx 0.76$ days.
\end{happeningbox}

\section{Classifying states: where can the chain go, and does it return?}
\label{ch23:classification}

Not all states are alike. Some you can leave and never come back to; some trap you; some you
visit forever. To organize this we ask two questions about every state: \emph{what is
reachable from it}, and \emph{if you start there, do you come back}?

\subsection{Communication and irreducibility}

\begin{definition}[Accessibility and communication]
\label{ch23:def-communicate}
State $j$ is \emph{accessible} from $i$, written $i \to j$, if $p_{ij}^{(n)} > 0$ for some
$n \ge 0$ -- you can get from $i$ to $j$ in some number of steps. States $i$ and $j$
\emph{communicate}, written $i \leftrightarrow j$, if $i \to j$ and $j \to i$.
\end{definition}

Communication is an equivalence relation (it is reflexive since $p_{ii}^{(0)}=1$, symmetric by
definition, and transitive by Chapman--Kolmogorov: if $i \to k$ in $r$ steps and $k \to j$ in
$s$ steps then $p_{ij}^{(r+s)} \ge p_{ik}^{(r)}p_{kj}^{(s)} > 0$). It therefore partitions the
state space into \emph{communicating classes} -- disjoint clusters of states within which
every member can reach every other.

\begin{definition}[Irreducible]
\label{ch23:def-irreducible}
A chain is \emph{irreducible} if all states communicate -- there is a single communicating
class, and you can get from anywhere to anywhere.
\end{definition}

On the state diagram, communicating classes are exactly the ``strongly connected components'':
the maximal blobs in which directed arrows let you travel both ways between any two members.

\subsection{Recurrent and transient states}

Now the return question. Start the chain at state $i$ and watch. Will it ever come back to
$i$? Let
\[
  f_i = \P(\text{chain returns to } i \text{ at some time} \ge 1 \given X_0 = i).
\]

\begin{definition}[Recurrent, transient]
\label{ch23:def-recurrent}
State $i$ is \emph{recurrent} if $f_i = 1$ (return is certain) and \emph{transient} if
$f_i < 1$ (there is positive probability you never return).
\end{definition}

\begin{intuitionbox}
A recurrent state is visited infinitely often with probability $1$, while a transient state is
visited only finitely many times before the chain departs from it permanently.
If the return probability is $f_i = 1$, then by the strong Markov property each subsequent
return is also certain, yielding infinitely many visits. If $f_i < 1$, the total number of
returns follows a geometric distribution with parameter $1 - f_i$, which is finite with
probability $1$.
\end{intuitionbox}

That intuition turns into a clean criterion. Let $N_i$ be the total number of visits to
$i$ (counting time $0$). Each time you are at $i$, the chance you ever return is $f_i$, and
returns are independent attempts because of the Markov property. So the number of returns is
geometric and
\[
  \E[N_i \given X_0 = i] = \sum_{n=0}^{\infty} p_{ii}^{(n)} = \frac{1}{1 - f_i},
\]
with the convention that the sum is $+\infty$ when $f_i = 1$. This gives a test you can run
straight from the matrix powers. In mathematical physics, a \emph{Green's function} (or potential kernel)
sums the total response of a system to a localized unit impulse. Here, the infinite series
$G(i,i) = \sum_{n=0}^{\infty} p_{ii}^{(n)}$ is the probabilistic Green's function: starting from
an initial unit mass at state $i$, it tallies the expected total time spent revisiting $i$
across the entire future.

\begin{theorem}[Recurrence via the Green's function]
\label{ch23:recurrence-test}
State $i$ is recurrent if and only if $\displaystyle\sum_{n=0}^{\infty} p_{ii}^{(n)} = \infty$,
and transient if and only if $\displaystyle\sum_{n=0}^{\infty} p_{ii}^{(n)} < \infty$.
Recurrence and transience are \emph{class properties}: if $i \leftrightarrow j$, then $i$ and
$j$ are either both recurrent or both transient.
\end{theorem}

\begin{whybox}
The expected number of visits is $\E[N_i] = \E\big[\sum_n \ind(X_n=i)\big] = \sum_n
\P(X_n=i \given X_0=i) = \sum_n p_{ii}^{(n)}$ by linearity (swap sum and expectation, allowed
for nonnegative terms). Separately, each visit independently leads to ``at least one more
visit'' with probability $f_i$, so $N_i$ is a geometric count with mean $1/(1-f_i)$. Setting
the two expressions equal: the sum is infinite exactly when $f_i = 1$. The class-property
claim follows because if $i \leftrightarrow j$ you can route visits to $i$ through $j$ and
back, so the two return sums are infinite together. The slogan: recurrence means you visit
infinitely often, and infinitely many visits means an infinite expected count.
\end{whybox}

For a \emph{finite} chain there is a satisfying shortcut: you cannot wander off to infinity,
because there is no infinity to wander to. With finitely many states and infinitely many time
steps, some state is visited infinitely often -- and a little argument shows every state in
its communicating class must be too.

\begin{proposition}[Finite chains: recurrence is about being trapped]
\label{ch23:finite-recurrence}
In a finite Markov chain, a communicating class is recurrent if and only if it is
\emph{closed} (no arrows lead out of it). At least one class is closed, so a finite chain
always has at least one recurrent state. Every finite irreducible chain is recurrent.
\end{proposition}

\begin{intuitionbox}
In a finite chain, recurrence is determined entirely by whether a communicating class has
transitions leaving it. If a class has transitions to other classes, the probability of
eventually leaving is $1$, and because the states cannot be re-entered from outside the class,
all states in that class are transient. If no transitions lead out of the class (the class is closed),
probability mass cannot leave, and every state in the class is recurrent.
\end{intuitionbox}

\begin{figure}[H]
\centering
\begin{tikzpicture}[->,>=Stealth,node distance=2.6cm,
  mystate/.style={circle,draw,thick,minimum size=0.95cm}]
  \node[mystate,draw=quotientorange] (1) {$1$};
  \node[mystate,draw=quotientorange,right=of 1] (2) {$2$};
  \node[mystate,draw=targetgreen,below=2.1cm of 1] (3) {$3$};
  \node[mystate,draw=targetgreen,below=2.1cm of 2] (4) {$4$};
  \path
    (1) edge[bend left=15] node[above]{} (2)
    (2) edge[bend left=15] node[below]{} (1)
    (1) edge node[left]{} (3)
    (2) edge node[right]{} (4)
    (3) edge[bend left=15] node[above]{} (4)
    (4) edge[bend left=15] node[below]{} (3);
  \node[quotientorange, above=7mm of $(1)!0.5!(2)$] {transient class $\{1,2\}$};
  \node[targetgreen, below=7mm of $(3)!0.5!(4)$] {recurrent class $\{3,4\}$};
\end{tikzpicture}
\caption{A four-state chain with two communicating classes. From $\{1,2\}$ (orange) there are
doors down into $\{3,4\}$ (green) but none coming back up, so $\{1,2\}$ is transient -- the
chain eventually falls through and never returns. The class $\{3,4\}$ is closed: no arrow
leaves it, so it is recurrent and the chain stays there forever after.}
\label{ch23:fig-classes}
\end{figure}

\begin{recoverybox}
\textbf{Pause and think: can all states in a chain be transient?}
Proposition~\ref{ch23:finite-recurrence} promises that every \emph{finite} chain possesses at least one recurrent state. Can an \emph{infinite} chain have every single state be transient?
\begin{itemize}
  \item Picture a particle hopping on the non-negative integers $\{0, 1, 2, \dots\}$ that always steps to the right: $p_{i, i+1} = 1$ for all $i \ge 0$.
  \item The particle marches strictly onward: $0 \to 1 \to 2 \to 3 \to \dots$. It visits each state $i$ exactly once (at time $n=i$) and never looks back. The return probability is $f_i = 0$ for every state!
  \item Every state is transient, and $\sum_{n=0}^\infty p_{ii}^{(n)} = p_{ii}^{(0)} = 1 < \infty$.
\end{itemize}
In a finite chain, the entire probability mass of $1$ must live on a finite set of nodes at every step, so with no escape route to infinity, the mass must eventually pool into at least one closed recurrent trap. On an infinite graph, mass can leak away to infinity forever.
\end{recoverybox}

\begin{example}[Gambler's ruin is mostly transient]
\label{ch23:ex-gamblers-ruin}
A gambler bets \$1 per round, winning with probability $p$ and losing with probability
$q = 1-p$, starting with $k$ and quitting at $0$ (broke) or $N$ (target). The state is the
current fortune, $S = \{0, 1, \dots, N\}$. The endpoints $0$ and $N$ are \emph{absorbing}:
$p_{00} = p_{NN} = 1$, each its own closed recurrent class. Every interior state
$1, \dots, N-1$ communicates with the others but has a door out (to $0$ or to $N$), so the
whole interior is \emph{transient}. The chain visits each fortune level only finitely often
before being absorbed. This is the structural reason the game must end: with probability $1$
the gambler is eventually ruined or hits the target. Computing \emph{which} is a beautiful
exercise in itself, left to the end of the chapter.
\end{example}

\begin{paradoxbox}[Polya's drunkard: recurrent in 2D, transient in 3D]
Consider a walker who starts at the origin and at each tick steps to a uniformly random
neighbor on the integer grid $\Z^d$. On the line ($\Z^1$) and in the plane ($\Z^2$) the
motion is \emph{recurrent}: with probability $1$ it returns to the origin (indeed to every point)
infinitely often. But in three dimensions ($\Z^3$) it is \emph{transient}: there is a positive
probability -- about $0.66$ -- that it \emph{never} comes home, and over time it drifts off to
infinity. George Polya proved this in 1921. Shizuo Kakutani summarized it perfectly: ``A drunk
man will find his way home, but a drunk bird may get lost forever.''

\textbf{Why the dimension matters.} Use the recurrence test
$\sum_n p_{00}^{(n)}$. The return probability at time $2n$ behaves like the chance a
spreading cloud of probability is back at its center. By the local central limit theorem the
walk after $2n$ steps is spread over a region of ``radius'' $\sqrt{n}$ in each dimension, so the
probability density at the origin scales like $n^{-d/2}$ in $d$ dimensions:
\[
  p_{00}^{(2n)} \;\asymp\; n^{-d/2}.
\]
Now sum: $\sum_n n^{-d/2}$ \emph{diverges} for $d = 1, 2$ (the exponents $1/2$ and $1$ are not
big enough) but \emph{converges} for $d \ge 3$ (exponent $3/2 > 1$). The whole dichotomy is the
convergence of a $p$-series. Two dimensions is exactly the borderline where the harmonic series
$\sum 1/n$ still diverges; one more dimension tips the sum into finiteness and lets the bird
escape.
\end{paradoxbox}

\begin{figure}[H]
\centering
\begin{tikzpicture}
\begin{axis}[
  width=12.5cm, height=5.6cm,
  xlabel={step $n$}, ylabel={position $X_n$},
  xmin=0, xmax=40, ymin=-6, ymax=7,
  xtick={0,10,20,30,40}, ytick={-6,-4,-2,0,2,4,6},
  title style={font=\small},
  title={Symmetric 1D random walk: recurrent paths returning to origin},
  grid=both, grid style={gray!18},
  every axis plot/.append style={thick}]
  \addplot[sourceblue] coordinates {
    (0,0)(1,-1)(2,-2)(3,-1)(4,0)(5,-1)(6,-2)(7,-3)(8,-2)(9,-1)(10,0)(11,-1)(12,0)(13,1)(14,2)(15,1)(16,2)(17,3)(18,2)(19,1)(20,0)(21,-1)(22,0)(23,1)(24,0)(25,-1)(26,-2)(27,-1)(28,-2)(29,-3)(30,-4)(31,-5)(32,-4)(33,-5)(34,-4)(35,-3)(36,-4)(37,-3)(38,-4)(39,-5)(40,-4)};
  \addplot[densityred] coordinates {
    (0,0)(1,1)(2,2)(3,1)(4,2)(5,3)(6,4)(7,5)(8,4)(9,5)(10,6)(11,5)(12,4)(13,5)(14,4)(15,3)(16,4)(17,3)(18,4)(19,3)(20,4)(21,3)(22,2)(23,1)(24,2)(25,1)(26,2)(27,3)(28,4)(29,3)(30,4)(31,5)(32,4)(33,3)(34,4)(35,5)(36,4)(37,5)(38,4)(39,3)(40,2)};
  \addplot[black,dashed,domain=0:40,samples=2]{0};
\end{axis}
\end{tikzpicture}
\caption{Two sample paths of the symmetric random walk on the line ($\Z^1$), one crossing the origin (dashed line) repeatedly and the other illustrating a long excursion away from it. In one and two dimensions every such path returns to its
start infinitely often with probability $1$ -- the recurrence of Polya's theorem made visible.
The vertical spread grows like $\sqrt{n}$, which is exactly the scaling behind the
$p_{00}^{(2n)} \asymp n^{-d/2}$ estimate in the box above.}
\label{ch23:fig-rw-paths}
\end{figure}

\subsection{Periodicity}

The last piece of structure is rhythm. Some states can only be revisited at multiples of a
fixed beat.

\begin{definition}[Period]
\label{ch23:def-period}
The \emph{period} of state $i$ is
\[
  d(i) = \gcd\set{ n \ge 1 : p_{ii}^{(n)} > 0 },
\]
the greatest common divisor of the times at which return is possible -- the
step-cadence that constrains returns. If $d(i) = 1$ the state
is \emph{aperiodic} (no rigid cadence); if $d(i) > 1$ it is \emph{periodic} with period $d(i)$.
Period is a class property: communicating states share a period.
\end{definition}

\begin{example}[The triangle walk has period... 2? No, 1]
\label{ch23:ex-triangle-period}
Recall the walk on a triangle that must move to a different corner each step (the Walking
question above). Return to corner $1$ is possible in $2$ steps ($1\to2\to1$) and in $3$ steps
($1\to2\to3\to1$). So the return-time set contains $2$ and $3$, and $\gcd(2,3) = 1$: the walk
is \emph{aperiodic}. Contrast a walk on a \emph{square} (a $4$-cycle) where you can only move
to an adjacent corner: returns to a corner happen only at even times, $\gcd = 2$, period $2$.
The triangle's odd loop breaks the rhythm; the square's bipartite structure enforces it
(Figure~\ref{ch23:fig-periodicity-triangle-square}).
\end{example}

\begin{figure}[H]
\centering
\begin{tikzpicture}[>=Stealth,
  mystate/.style={circle,draw=#1,thick,fill=#1!12,minimum size=0.95cm,font=\small\bfseries}]
  % Left panel: Triangle (Aperiodic)
  \begin{scope}[shift={(0,0)}]
    \node[mystate=sourceblue] (T1) at (0, 2.2) {$1$};
    \node[mystate=sourceblue] (T2) at (-1.8, -0.6) {$2$};
    \node[mystate=sourceblue] (T3) at (1.8, -0.6) {$3$};

    \draw[<->, thick, sourceblue] (T1) -- node[above left=1pt, font=\scriptsize] {$1/2$} (T2);
    \draw[<->, thick, sourceblue] (T2) -- node[below=2pt, font=\scriptsize] {$1/2$} (T3);
    \draw[<->, thick, sourceblue] (T3) -- node[above right=1pt, font=\scriptsize] {$1/2$} (T1);

    \node[font=\small\bfseries, align=center] at (0, -1.5) {Triangle: Aperiodic ($d=1$)};
    \node[font=\footnotesize, gray!40!black, align=center] at (0, -2.2) {
      Return times: $2$ ($1 \leftrightarrow 2$), $3$ ($1 \to 2 \to 3 \to 1$)\\
      $\gcd(2, 3) = 1 \implies$ rhythm broken
    };
  \end{scope}

  % Divider between panels
  \draw[dashed, gray!40, thick] (3.4, -2.5) -- (3.4, 2.6);

  % Right panel: Bipartite Square (Period 2)
  \begin{scope}[shift={(6.8,0)}]
    \node[mystate=sourceblue] (S1) at (-1.4, 1.4) {$1$};
    \node[mystate=quotientorange] (S2) at (1.4, 1.4) {$2$};
    \node[mystate=sourceblue] (S3) at (1.4, -0.6) {$3$};
    \node[mystate=quotientorange] (S4) at (-1.4, -0.6) {$4$};

    \draw[<->, thick, gray!60!black] (S1) -- node[above=2pt, font=\scriptsize] {$1/2$} (S2);
    \draw[<->, thick, gray!60!black] (S2) -- node[right=2pt, font=\scriptsize] {$1/2$} (S3);
    \draw[<->, thick, gray!60!black] (S3) -- node[below=2pt, font=\scriptsize] {$1/2$} (S4);
    \draw[<->, thick, gray!60!black] (S4) -- node[left=2pt, font=\scriptsize] {$1/2$} (S1);

    \node[font=\small\bfseries, align=center] at (0, -1.5) {Square: Periodic ($d=2$)};
    \node[font=\footnotesize, gray!40!black, align=center] at (0, -2.2) {
      Bipartite: blue $\leftrightarrow$ orange\\
      Return times: $\{2, 4, 6, \dots\}$, $\gcd=2$
    };
  \end{scope}
\end{tikzpicture}
\caption{Periodicity is about odd versus even cycles. Left: on a triangle, a particle can return to state $1$ in $2$ steps ($1 \to 2 \to 1$) or in $3$ steps ($1 \to 2 \to 3 \to 1$), giving $\gcd(2,3) = 1$; the odd loop breaks any rigid beat, making the chain aperiodic. Right: on a bipartite square, every step alternates between blue (odd parity) and orange (even parity). Returning to state $1$ requires an even number of steps, giving $\gcd(2,4,6,\dots) = 2$ and locking the distribution into a persistent two-step alternation.}
\label{ch23:fig-periodicity-triangle-square}
\end{figure}

\begin{happeningbox}
Period is really a statement about the eigenvalues lying \emph{on} the unit circle. An
aperiodic irreducible chain has $\lambda = 1$ as its only eigenvalue of modulus $1$, so
$P^n$ settles down to a single limiting matrix. A chain with period $d$ has $d$ eigenvalues
spread evenly around the unit circle (the $d$-th roots of unity), and $P^n$ never settles --
it cycles through $d$ phases forever. The bipartite square chain has eigenvalues $\pm 1$; the
$+1$ gives the equilibrium and the $-1$ gives the persistent odd/even alternation. Aperiodicity
rules out these unit-circle oscillations and lets the chain converge, which
is why the next chapter's convergence theorem assumes it.
\end{happeningbox}

\begin{aestheticbox}
The structural properties of a Markov chain organize as follows:
\begin{itemize}
  \item \textbf{Communicating classes}: the strongly connected components of the state diagram.
    A chain is \emph{irreducible} if there is only one class.
  \item \textbf{Closed vs.\ open}: in a finite chain, a class is recurrent if and only if no arrow
    leaves it; open classes are transient.
  \item \textbf{Period}: the greatest common divisor of possible return times; \emph{aperiodic}
    means period $1$.
\end{itemize}
A finite, irreducible, aperiodic chain is called \emph{ergodic}. These chains converge to a unique
stationary equilibrium from any initial distribution: irreducibility guarantees all states communicate,
recurrence ensures no probability mass escapes to transient states, and aperiodicity rules out
rigid cyclic oscillations. Each property can be verified directly from the state diagram.
\end{aestheticbox}

\section{A larger worked example}
\label{ch23:worked-example}

Time to put everything to work on the reader-mood chain of Example~\ref{ch23:ex-skim},
\[
  P = \begin{pmatrix}
   0.7 & 0.2 & 0.1 \\
   0.3 & 0.5 & 0.2 \\
   0.2 & 0.3 & 0.5
  \end{pmatrix},
  \qquad \text{states } F, K, D.
\]

\textbf{Step 1: classify.} Every entry of $P$ is positive, so $i \to j$ in one step for all
$i,j$: the chain is irreducible (one class) and, since $p_{ii}>0$ gives return in one step,
aperiodic (period $\gcd$ includes $1$). It is finite and irreducible, hence recurrent. This is
an ergodic chain -- it will settle to a unique equilibrium.

\textbf{Step 2: two- and three-step behavior.} Suppose you start Focused, $\mu^{(0)} =
(1,0,0)$. Then
\[
  \mu^{(1)} = (0.7,\,0.2,\,0.1), \qquad
  \mu^{(2)} = \mu^{(1)}P = (0.57,\,0.27,\,0.16),
\]
\[
  \mu^{(3)} = \mu^{(2)}P = (0.512,\,0.297,\,0.191).
\]
(For instance $\mu^{(2)}_F = 0.7\cdot0.7 + 0.2\cdot0.3 + 0.1\cdot0.2 = 0.49+0.06+0.02 =
0.57$.) The probability of being focused is sliding down from $1$ toward its long-run value.

\textbf{Step 3: the equilibrium.} The long-run distribution $\pi = (\pi_F, \pi_K, \pi_D)$ is
the one that does not change under a step: $\pi P = \pi$ with $\pi_F + \pi_K + \pi_D = 1$.
Writing out $\pi P = \pi$:
\[
  0.7\pi_F + 0.3\pi_K + 0.2\pi_D = \pi_F, \quad
  0.2\pi_F + 0.5\pi_K + 0.3\pi_D = \pi_K, \quad
  0.1\pi_F + 0.2\pi_K + 0.5\pi_D = \pi_D.
\]
Solving (any two equations plus normalization) gives the exact answer
\[
  \pi = \Big(\tfrac{19}{41},\, \tfrac{13}{41},\, \tfrac{9}{41}\Big)
      \approx (0.463,\, 0.317,\, 0.220).
\]
In the long run you spend about $46.3\%$ of paragraphs focused, $31.7\%$ skimming, $22.0\%$
distracted -- regardless of how you began. Compare with $\mu^{(3)} = (0.512, 0.297, 0.191)$:
after just three paragraphs the distribution is already in the right neighborhood. The figure
shows the full approach.

\begin{figure}[H]
\centering
\begin{tikzpicture}
\begin{axis}[
  width=12.5cm, height=6.4cm,
  xlabel={paragraph $n$}, ylabel={probability of each mood},
  xmin=0, xmax=12, ymin=0, ymax=1.05,
  xtick={0,2,4,6,8,10,12},
  title style={font=\small},
  title={Reading mood chain: relaxation to stationary equilibrium $\pi$},
  legend pos=north east, grid=both, grid style={gray!18},
  every axis plot/.append style={thick,mark size=1.3pt}]
  \addplot[sourceblue,mark=*] coordinates {
    (0,1.0000)(1,0.7000)(2,0.5700)(3,0.5120)(4,0.4857)(5,0.4737)(6,0.4681)(7,0.4656)(8,0.4644)(9,0.4639)(10,0.4636)(11,0.4635)(12,0.4635)};
  \addlegendentry{Focused}
  \addplot[quotientorange,mark=square*] coordinates {
    (0,0.0000)(1,0.2000)(2,0.2700)(3,0.2970)(4,0.3082)(5,0.3131)(6,0.3152)(7,0.3162)(8,0.3167)(9,0.3169)(10,0.3170)(11,0.3170)(12,0.3171)};
  \addlegendentry{Skimming}
  \addplot[densityred,mark=triangle*] coordinates {
    (0,0.0000)(1,0.1000)(2,0.1600)(3,0.1910)(4,0.2061)(5,0.2133)(6,0.2166)(7,0.2182)(8,0.2189)(9,0.2192)(10,0.2194)(11,0.2195)(12,0.2195)};
  \addlegendentry{Distracted}
  \addplot[sourceblue,dashed,domain=0:12,samples=2]{0.4634};
  \addplot[quotientorange,dashed,domain=0:12,samples=2]{0.3171};
  \addplot[densityred,dashed,domain=0:12,samples=2]{0.2195};
\end{axis}
\end{tikzpicture}
\caption{Starting Focused (the blue curve begins at probability $1$), the distribution
$\mu^{(n)} = \mu^{(0)}P^n$ relaxes to the equilibrium $\pi = (19/41, 13/41, 9/41) \approx
(0.463, 0.317, 0.220)$ (dashed lines). Because the chain is ergodic, the same three dashed limits are reached from
\emph{any} starting mood; only the transient first few paragraphs differ.}
\label{ch23:fig-reader-converge}
\end{figure}

\begin{applicationbox}[PageRank: a Markov chain that ranks the web]
Google's original PageRank algorithm is one giant Markov chain. The states are web pages; the
transitions model a ``random surfer'' who, from a page, clicks a uniformly random outgoing
link. The transition matrix $P$ has $p_{ij} = 1/(\text{number of links out of } i)$ for each
link $i \to j$. With billions of pages this matrix is astronomically large, but it is sparse
(most entries zero) and we never form $P^n$ explicitly -- we just iterate $\mu^{(n+1)} =
\mu^{(n)}P$, which is one sparse matrix-vector multiply, the operation of
Section~\ref{ch23:matrix-powers}.

Two problems arise, and both are fixed with the classification language of this chapter. Pages
with no outgoing links are absorbing dead-ends (a closed class of size one) that would soak up
all the probability; and the link graph is not irreducible, so a naive surfer can get trapped in
a corner of the web. Both are cured by \emph{teleportation}: with probability $1 - \alpha$
(typically $1-\alpha = 0.15$) the surfer jumps to a uniformly random page instead of following a
link. The modified matrix
\[
  G = \alpha P + (1-\alpha)\tfrac{1}{m}\mathbf{1}\mathbf{1}^{\!\top}
\]
is strictly positive, hence irreducible and aperiodic -- ergodic. Its unique equilibrium $\pi$
exists, is unique, and is reached fast (the teleportation makes the second eigenvalue at most
$\alpha = 0.85$, so the iteration converges geometrically). \emph{That equilibrium vector
$\pi$ is the PageRank}: $\pi_i$ is the long-run fraction of time the random surfer spends on
page $i$, and pages where the surfer lingers are ranked higher. A multibillion-dollar idea is,
at bottom, the stationary distribution of an ergodic Markov chain.
\end{applicationbox}

\begin{connectionbox}
In the previous chapter on modes of convergence, we studied what it means for random variables
to converge in distribution: the probabilities $\P(X_n = j)$ march toward a fixed limit $\pi_j$.
In this chapter, we saw a concrete instance of that convergence: the state distributions
$\mu^{(n)} = \mu^{(0)} P^n$ settle toward a unique equilibrium vector $\pi$.

That equilibrium -- finding the row vector $\pi$ with $\pi P = \pi$, proving it exists and is
unique, and calculating how rapidly $\mu^{(n)} \to \pi$ -- is the central theme of the next
chapter, \emph{Markov Chain Steady States and Mixing}. There we will explore:
\begin{itemize}
  \item \textbf{The convergence theorem}: why every finite, irreducible, aperiodic chain
    forgets its past and settles into its stationary distribution.
  \item \textbf{Mixing times and spectral gaps}: how the second eigenvalue $\lambda_2$ governs the
    clock of convergence, and how cutoff phenomena create abrupt transitions to equilibrium.
  \item \textbf{Coupling}: a probabilistic technique that proves convergence by running
    two copies of the chain on the same graph until their paths meet, bypassing matrix powers directly.
\end{itemize}
Everything there stands on the three pillars established here: path probabilities multiply
(Proposition~\ref{ch23:path-factorization}), $n$-step transitions are matrix powers
(Theorem~\ref{ch23:chapman-kolmogorov}), and the classification triad (irreducible, recurrent,
aperiodic) guarantees that the equilibrium exists. Keep the state diagram in your mind's eye;
it is the roadmap for all that follows.
\end{connectionbox}

\section{Exercises}
\label{ch23:exercises}

\begin{exercise}[Read a matrix]
\label{ch23:ex-read-matrix}
A mouse moves among three rooms $A, B, C$ with
$P = \left(\begin{smallmatrix} 0 & 1 & 0 \\ 1/2 & 0 & 1/2 \\ 0 & 1 & 0 \end{smallmatrix}\right)$.
(a) Draw the state diagram. (b) Verify $P$ is stochastic. (c) Is the chain irreducible?
(d) Compute the period of state $A$ and explain why this chain never settles to a single
distribution.
\end{exercise}

\begin{exercise}[Two-step probabilities by hand]
\label{ch23:ex-twostep}
For the weather chain $P = \left(\begin{smallmatrix} 0.8 & 0.2 \\ 0.4 & 0.6 \end{smallmatrix}\right)$,
compute $P^3$ and find $\P(X_3 = \text{Rain} \given X_0 = \text{Sun})$. Then verify your
answer for $p_{SS}^{(3)}$ by listing every length-$3$ path $S \to \cdot \to \cdot \to S$ and
summing their probabilities.
\end{exercise}

\begin{exercise}[Push a distribution]
\label{ch23:ex-push}
Using the reader-mood matrix $P$ of Section~\ref{ch23:worked-example} and the starting
distribution $\mu^{(0)} = (0,0,1)$ (you begin Distracted), compute $\mu^{(1)}$ and
$\mu^{(2)}$. Does your $\mu^{(2)}$ look like it is heading to the same $\pi = (19/41,
13/41, 9/41) \approx (0.463, 0.317, 0.220)$ as the Focused start? Explain why it must, in one sentence.
\end{exercise}

\begin{exercise}[Memoryless, carefully]
\label{ch23:ex-memoryless}
Let $(X_n)$ be a Markov chain. (a) Argue from Definition~\ref{ch23:def-markov} that
$\P(X_2 = k \given X_0 = i, X_1 = j) = \P(X_2 = k \given X_1 = j)$. (b) Give a concrete chain
and states for which $\P(X_2 = k \given X_0 = i) \ne \P(X_2 = k)$, demonstrating that the
past and future are \emph{not} independent without conditioning on the present.
\end{exercise}

\begin{exercise}[Classify a chain]
\label{ch23:ex-classify}
A chain on $\{1,2,3,4,5\}$ has transitions: $1 \to 2$ (prob $1$); $2 \to 1$ or $2 \to 3$
(each $1/2$); $3 \to 3$ (prob $1$); $4 \to 5$ and $5 \to 4$ (each prob $1$). (a) Find all
communicating classes. (b) Which are closed? (c) Which states are recurrent and which
transient? (d) Find the period of each recurrent class.
\end{exercise}

\begin{exercise}[Gambler's ruin]
\label{ch23:ex-ruin}
In the gambler's ruin chain of Example~\ref{ch23:ex-gamblers-ruin} with fair bets
($p = q = 1/2$), let $h_k$ be the probability of reaching the target $N$ before going broke,
starting from fortune $k$. (a) Argue by conditioning on the first bet that
$h_k = \tfrac12 h_{k-1} + \tfrac12 h_{k+1}$ for $1 \le k \le N-1$, with $h_0 = 0$ and
$h_N = 1$. (b) Solve this recurrence to get $h_k = k/N$. (c) Interpret: why is the chance of
success just your fortune as a fraction of the target?
\end{exercise}

\begin{exercise}[Recurrence sum]
\label{ch23:ex-recurrence-sum}
A chain on $\{1, 2, 3, \dots\}$ has $p_{i,i+1} = 1$ for all $i$ (always step right, never
return). (a) Compute $\sum_{n=0}^\infty p_{11}^{(n)}$. (b) Using
Theorem~\ref{ch23:recurrence-test}, conclude whether state $1$ is recurrent or transient.
(c) Reconcile this with Proposition~\ref{ch23:finite-recurrence} -- why does the ``finite
chains always have a recurrent state'' guarantee not apply here?
\end{exercise}

\begin{exercise}[Eigenvalues and forgetting]
\label{ch23:ex-eigen}
For a general two-state chain $P = \left(\begin{smallmatrix} 1-a & a \\ b & 1-b
\end{smallmatrix}\right)$ with $0 < a, b < 1$: (a) Verify that $\mathbf{1} = (1,1)^\top$ is a
right eigenvector with eigenvalue $1$. (b) Show the other eigenvalue is $\lambda_2 = 1 - a -
b$ and note $\abs{\lambda_2} < 1$. (c) Explain, using
$P^n = V \Lambda^n V^{-1}$, why the chain forgets its start at rate $\abs{1-a-b}^n$, and what
happens in the special case $a + b = 1$.
\end{exercise}

\begin{exercise}[Period by gcd]
\label{ch23:ex-period-gcd}
A chain has $p_{11}^{(n)} > 0$ exactly when $n \in \{4, 6, 8, 10, \dots\}$ (every even number
$\ge 4$). (a) What is the period of state $1$? (b) Now suppose instead returns are possible at
$\{6, 10, 15\}$ and their sums; what is the period? (c) Conceptually, why does adding a single
odd return time of any length force aperiodicity once it shares no common factor with the
rest?
\end{exercise}

\begin{exercise}[Build a model, prove a property]
\label{ch23:ex-build}
The Ehrenfest model of diffusion: $N$ molecules are split between two chambers; each second,
one molecule chosen uniformly at random switches chambers. Let $X_n$ be the number in the
left chamber, a chain on $\{0, 1, \dots, N\}$. (a) Write the transition probabilities
$p_{i,i-1}$ and $p_{i,i+1}$. (b) Show the chain is irreducible. (c) Show it has period $2$, and
explain physically what the period means. (d) \emph{(Proof.)} Show that the binomial
distribution $\pi_i = \binom{N}{i} 2^{-N}$ satisfies the \emph{detailed balance} condition
$\pi_i p_{i,i+1} = \pi_{i+1} p_{i+1,i}$ -- which says the probability flux flowing forward
from $i$ to $i+1$ exactly balances the reverse flux from $i+1$ to $i$ -- and explain why summing
these two-way fluxes across all states makes $\pi$ stationary ($\pi P = \pi$).
\end{exercise}
